%!TEX root = ../../main.tex
\section{기여와 논문의 구성}
\label{sec:intro-contrib}

본 논문의 기여는 다음 네 가지이다.

\begin{enumerate}[leftmargin=2em]
  \item \textbf{근거 있는 교전 사례 구성 (M-RQ1; 제 \ref{sec:eng-temporal}--\ref{sec:eng-scale} 절, \ref{sec:eng-sensitivity} 절).} 시간 경계 $G$와 공간 경계 $D$를 코퍼스에서 추정하고, 게이트 상수 $R$, $B$를 게임 규칙에 고정하며, 교전 집합을 결정하는 모든 상수를 출처 분류와 함께 표로 남긴다. 규모 계급을 작은 쪽 참여 인원으로 정의하여 한타 T와 소규모 교전 S를 별개의 코호트로 두고, 정의 상수의 일변량 변경이 사례 구성을 어떻게 바꾸는지 센다.
  \item \textbf{경기 승패와 연결된 결과 가치 (M-RQ2; 제 \ref{sec:value-evaluator}--\ref{sec:value-shold} 절).} 경기 승패로 학습해 동결한 승률 평가기 $\Vhat$로 교전 구간 전후의 추정 승률 변화 방향 $\Ysvi = \ind[\dV>0]$을 정의하고, 평가기 품질, 변화의 방향·평균·규모의 분리, 비교전 구간 대조, 물질적 결과와 다른 결과 정의와의 대응, 지평·평가기·종점에 대한 라벨의 안정성, 시간 경과와 관측 갱신의 분해로 이 측정이 무엇을 재는지 한정한다.
  \item \textbf{초기 우세와 시간을 넘는 예측 효용 (M-RQ3; 제 \ref{sec:pred-primary-T}--\ref{sec:pred-infogroups} 절).} 코호트별로 교전 전 상태의 정규화 로지스틱 모델 $q$가 승률·시간의 스플라인 기준선 $\PTflex$보다 Brier 점수를 낮추는지 동일 행의 짝 대비와 경기 단위 부트스트랩 구간으로 답하고, 한타 모델의 전이와 합동 학습의 효과를 이차 대비로 보고한다. 동일한 학습 절차에서 프레임 기반 특징과 사건 기반 특징의 추가 예측 효용을 비교한다. 학습기 간 정합 비교는 후속 검증 항목으로 남아 있다.
  \item \textbf{의미와 범위의 검증 (M-RQ4; 제 \ref{sec:val-b40}--\ref{sec:val-external} 절).} 균형 구간, 작은 변화 제외, 관측 갱신 층화, 점수 분해, 이후 패치·다른 지역의 점수 전용 평가와 그 부트스트랩 구간으로 이득이 어디까지 유지되는지 보고한다.
\end{enumerate}

각 결과의 해석 범위는 제 \ref{sec:disc-scope} 절에 한 곳에 둔다.

논문의 구성은 다음과 같다. 제 \ref{ch:related} 장은 관련 연구를 승률 모형, 교전 검출과 교전 단위 예측, 사건 가치, 적정 점수와 짝 비교의 네 주제로 정리하고 남은 문제와 본 논문의 위치를 밝힌다. 제 \ref{ch:data} 장은 비동기 공개 텔레메트리의 구조, 원천 코퍼스, 데이터 역할, 교전 전 입력의 구성을 다룬다. 제 \ref{ch:engagement} 장은 교전 사례의 정의, 상수의 출처, 규모 계급과 코호트, 결과 시점 규칙, 분석 대상, 정의 민감도를 다룬다. 제 \ref{ch:value} 장은 동결 평가기와 $\SVI$의 정의, 그리고 그 측정이 무엇을 재는지에 대한 검증을 제시한다. 제 \ref{ch:prediction} 장은 예측기, 기준선, 선정·보정 절차, 코호트별 주 결과, 전이·합동 대비, 정보군 비교를 제시한다. 제 \ref{ch:validation} 장은 균형 구간, 작은 변화 제외, 관측 갱신 층화, 점수 분해, 외부 평가로 예측의 의미와 범위를 검증한다. 제 \ref{ch:discussion} 장은 결과의 연결, 해석의 범위, 평가기 의존성, 한계, 향후 과제를 논의하고 결론을 맺는다. 부록에는 재현 자료와 무결성 검사, 상수표, 데이터 역할 census, 수치의 증거 추적, 실행 상태를 둔다.
